% DRAFT: Q4 canonical 已冻结；Q1/Q2/Q3 精确结果及 Q4 正式场景需在各自 registry FROZEN+CHECKED 后统一补全。
本文针对六区域人工智能算力与电力系统协同调度问题，首先从附件中识别任务、区域与能源三类非均匀结构：AITraining 的任务数与真实计算工作量明显错位，六区域的算力、网络时延、PUE 与储能/售电优势并不重合，且基准状态存在大量“弃电与供负荷购电并存”的时空错配。由此，本文不采用主观多指标加权，而按照“需求结构识别--计算柔性--储能柔性--算储联合”的递进路线建立模型。

\textbf{针对问题一，}根据小时到达近齐次 Poisson、TaskType 与 SourceRegion 显著关联以及任务数与 GPU-hour 错位的统计特征，建立联合业务标记与条件资源标记的复合 Poisson 短期预测模型，并对 2376--2399\,h 的真实任务实施基础算力调度。模型同时输出层级一致的 24\,h 概率区间、任务级调度方案、最后 24\,h 甘特图和区域 GPU 利用率，为后续能源感知调度提供 workload baseline。

\textbf{针对问题二，}在储能与既有能源动作固定条件下释放 workload 的跨区与开工时段自由度，利用分段动态边际能源响应和完整合法域迭代扫描构造碳感知调度。Cost-primary 与 Carbon-primary 两个端点总体同向但存在可测边际冲突；多起点结果同时表明精确启发式端点具有路径依赖，因此本文将其作为 workload-only 高质量可行基线，而不包装为全局整数最优。

\textbf{针对问题三，}固定任务负荷后建立含新能源直接供电、储能充放电、购售电与 SOC 约束的线性能源调度模型。由新能源盈余 $H_{rt}$ 与 SellLimit 的数据结构自然识别 A/B/C、D/E、F 三类储能角色，并通过 B0/E0/E1 三层对照将 raw reference、无储能能源优化和 BESS 增量价值严格分离。

\textbf{针对问题四，}由于 workload 迁移与 BESS 充电会争用同一时空新能源余量，建立统一的 workload--energy--storage 联合模型，并采用 region multi-cut Benders 与 exact column generation 隐式搜索约 $2.33\times10^8$ 个合法 placements，再按 Cost$\rightarrow$Wait$\rightarrow$Latency 进行字典序精化。50,000 任务的 final-pool 两轮再认证得到真实能源成本 $-4.5934\times10^8$ CNY、Total Wait=0、总网络时延 250,860 ms，仅 61 个任务跨区执行（0.122\%）。最终结果显示，\textbf{储能时间柔性替代了绝大部分原本由 workload 时间平移与跨区迁移承担的调节需求}。完整域 Latency LP 下界与整数代表解的相对差约 0.156\%，故本文不声称 50k 全局整数最优。

\textbf{本文特点：}以附件数据关系决定模型结构，不对 Cost、Carbon、QoS、Latency、新能源利用率和峰值购电进行主观加权；Q2/Q3 分别识别计算柔性与储能柔性，Q4 再检验两者的耦合与替代；同时通过独立预测检验、多起点审计、E0/E1 对照、40-task exact benchmark、完整域定价和真实 Energy recourse 复核保证结论可追溯。Q4 的碳约束、电价机制与新能源波动正式场景将在同一联合模型下完成后统一冻结。

\keywords{算电协同；异构算力调度；储能优化；Benders分解；列生成；字典序优化}